% !TEX root = ../main.tex
\section{Discussion}
\label{sec:discussion}

\subsection{Limitations}
\label{sec:limitations}

\paragraph{Single-agent mutation.}
Each experiment cycle invokes exactly one mutation agent that edits the artifact in isolation, so mutations that require coordinated changes across interacting files cannot be generated atomically: an edit to a function signature and its call sites must either be proposed as a single multi-file patch by one agent or deferred across multiple experiment cycles.
The current architecture does not support multi-agent coordination with conflict resolution or merge semantics, so optimization targets whose improvements fundamentally require parallel specialization (one agent refactoring structure while another refines documentation) fall outside the framework's reach.

\paragraph{Text artifacts only.}
The engine operates on file-based text artifacts versioned in git.
Binary files (images, compiled assets, database schemas), structured data stores, and deployed services are outside the mutation scope.
Optimization targets requiring deployment to external systems (email campaigns, ad platforms) are supported only as read-only proposal generation with approval gates; this pathway is implemented but not evaluated here, and it is not integrated with external deployment systems.

\paragraph{LLM-as-judge evaluation quality.}
Stochastic evaluation relies on an LLM judge scoring against binary criteria.
The quality of optimization is bounded by the quality of the judge: a judge that rewards superficial keyword matching will drive the optimizer toward keyword-stuffed artifacts.
Position debiasing (\cref{sec:eval-pipeline}) and held-out validation mitigate systematic judge bias, but they do not eliminate it.
Practitioners must invest in evaluation design (clear, unambiguous criteria with concrete acceptance thresholds) to obtain meaningful optimization trajectories.

\paragraph{Cost.}
A full optimization run of 50 experiments incurs substantial LLM API cost, with the exact figure depending on the evaluation mode, the number of stochastic samples per experiment, and the selected models for the mutation, diagnosis, and judgment roles.
Two-phase mutation and multi-draft generation can increase per-experiment cost relative to single-phase greedy mutation because of the additional LLM calls for diagnosis and draft generation.
For cost-sensitive applications, the engine supports cheaper diagnosis models (e.g., a smaller model for diagnosis and a larger one for mutation) and adaptive sample sizing that reduces evaluation cost for clear winners and losers.

\paragraph{Cold start.}
The knowledge system provides no benefit during the first several experiments, when the experiment history is empty.
The TF-IDF retrieval index requires a minimum corpus size before similarity queries become discriminative, and the engine's cold-start protocol injects only raw experiment records until a configurable activation threshold is reached.

\subsection{Threats to Validity}
\label{sec:threats}

\paragraph{Internal validity.}
LLM outputs are nondeterministic: the same prompt produces different mutations across runs.
We control for this by pairing runs by random seed across configurations and using nonparametric statistical tests that do not assume distributional form.
API pricing variability, model availability changes, backend switching, rate limits, and transient retry behavior introduce noise into cost measurements.
The archived end-to-end records include mixed model routing (\cref{tab:result-provenance}), so the cost tables should be interpreted as telemetry for that result bundle rather than as provider-independent cost estimates.
The analysis mitigates billing opacity by reading recorded per-experiment cost fields instead of relying on billing-console totals.

\paragraph{Empirical scope and statistical power.}
The end-to-end comparison in \cref{sec:end-to-end} runs at $n = 5$ paired seeds per target across five targets.
At this seed count the Wilcoxon signed-rank test has a two-sided $p$-value floor of $0.0625$, so no comparison can reach the conventional $\alpha = 0.05$ threshold per-target without ties or unusual sign patterns, and after Holm--Bonferroni correction across five targets the bar rises further still.
The directional pattern across targets (treatment beats control on 5/5 targets, treatment beats greedy on 5/5 targets) and the Cohen's-large effect sizes on the two strongest greedy-versus-treatment comparisons (B3, $d = -1.09$; B5, $d = 1.35$) indicate that seed count is one binding constraint on whether systemic claims clear the corrected threshold.
A higher-powered replication (larger $n$ per target, or a study restricted to one or two targets that does not pay the five-target multiplicity tax) would test whether the directional advantage of the treatment configuration translates into a statistically defensible claim.
We have intentionally not run such a replication at this stage because the seed-counted result with its honest non-significance is more informative than a cherry-picked single-target study would be: it illustrates exactly the conservatism the framework's statistical-acceptance layer is designed to enforce.
The domain-agnostic design claim in \cref{sec:introduction} continues to rest primarily on the architecture of the artifact--evaluation--agent triplet (\cref{sec:triplet}); the suite spans prompts, code, and documentation, but five targets cannot establish that the framework composes into improvements across the broader space of text-artifact optimization problems.

\paragraph{External validity.}
The five benchmark targets that ship with the framework are themselves a simplification of real-world optimization problems.
Production artifacts are larger, have more complex dependency structures, and operate under constraints (backward compatibility, style guides, regulatory requirements) that the benchmark targets do not capture.
The benchmark suite is intended as a reproducible reference workload for comparing configurations under controlled conditions, not as a representative sample of the distribution of problems a practitioner would encounter in practice.
The suite is useful for controlled comparison and failure analysis, but it does not establish that relative ordering of configurations will transfer unchanged to production-scale artifacts.

\paragraph{Construct validity.}
Stochastic evaluation quality depends on the specificity and completeness of the scoring criteria.
Vague criteria (e.g., ``is the output high quality?'') produce noisy scores that obscure real improvements.
The benchmark targets ship with criteria that are phrased as concrete yes/no questions with explicit acceptance conditions, and the position-bias cancellation and Bradley--Terry scoring described in \cref{sec:eval-pipeline} further reduce judge-level artefacts.
This investment in criterion design is itself a limitation: in practice, users may not invest equivalent effort, leading to lower optimization effectiveness than the benchmark design suggests.

\paragraph{Proxy gaming and rubric-facing optimization.}
The architectural separation in \cref{sec:triplet} (immutable eval definitions, an agent that never sees the eval command, scope enforcement that blocks edits to scoring code) prevents the agent from modifying the evaluation function itself.
It does not prevent the agent from optimizing the evaluation function as it is written.
On stochastic-eval targets, an agent can learn to produce artifacts that score well on the rubric criteria as phrased even when the produced artifact does not match the rubric's intent: rubric-facing language that satisfies a binary keyword check, formatting choices that satisfy a structural criterion without addressing the underlying quality the criterion is meant to proxy, or coverage of all listed cases without depth on any of them.
On deterministic-eval targets, an agent can overfit to the test fixtures: passing the specific inputs the harness exercises while degrading on inputs the harness does not exercise.
The framework's defense against these failure modes is external to the optimization loop---held-out divergence detection (\cref{sec:eval-pipeline}) flags artifacts whose held-out and primary-set scores diverge---but the user remains responsible for designing rubrics and harnesses where the proxy is closely aligned with the intent.
Where alignment is weak, the framework will optimize the proxy, not the intent.

\subsection{Future Work}
\label{sec:future-work}

\paragraph{Higher-powered empirical replication.}
The seed-counted end-to-end study in \cref{sec:end-to-end} establishes a consistent directional advantage of the treatment configuration over both control and greedy across the suite, but does not clear the multiplicity-corrected significance threshold at $n = 5$.
The two strongest greedy-versus-treatment comparisons (B3, B5) sit at the Wilcoxon two-sided $p$-value floor for $n = 5$, indicating that statistical power is one unresolved constraint on the systemic claim.
A replication at higher seed counts on a smaller subset of targets would resolve whether the directional advantage translates into a statistically defensible claim.

\paragraph{Multi-artifact optimization.}
Extending the mutation scope to coordinated changes across multiple files (e.g., modifying a function signature and all its call sites simultaneously) would enable structural refactoring that single-file mutations cannot achieve.
This requires dependency-aware scoping and atomic multi-file commits with rollback semantics, both of which the git environment backend supports in principle.

\paragraph{Typed evaluator protocol.}
The current evaluation interface is command-based: the engine invokes a shell command and parses its stdout.
A typed Python callable protocol (already prototyped in the codebase) would enable programmatic evaluation with structured return types, eliminating parsing fragility and supporting evaluation functions that require in-process state (database connections, loaded ML models).

\paragraph{MAP-Elites integration.}
The quality-diversity archive (\cref{sec:enhancements}) maintains a grid of behaviorally diverse solutions but is not yet integrated into the search loop as a solution source.
Full MAP-Elites integration would allow the engine to optimize for diversity and quality simultaneously, producing a portfolio of distinct high-scoring artifacts rather than a single best solution.

\paragraph{Continuous optimization.}
A daemon mode that monitors artifact changes and triggers optimization runs on detected regressions would enable always-on quality maintenance.
The scheduler and daemon management components are implemented; the remaining work is policy design for when to trigger optimization (commit hooks, schedule, metric drift detection) and how to surface results without disrupting the developer workflow.

\paragraph{IDE integration.}
Embedding the optimization loop into editor extensions (VS Code, JetBrains) would reduce the barrier to adoption by allowing developers to trigger optimization on selected files without leaving their editor.
The engine's file-based architecture (JSONL logs, file-based SSE dashboard) is compatible with extension integration, but the UX design (progress visualization, result acceptance, inline diff preview) remains open.
